\chapter{Continuum and Kinetic Formulation}
\label{chap:theory}

\section{The Boussinesq problem}
\label{sec:theory-boussinesq}

Let $\mathbf{u}(\mathbf{x},t)$ be the velocity field, $p$ the kinematic
pressure, and $\theta$ the buoyancy scalar measured as a temperature deviation
from a linear background profile $\bar{\theta}(z) = \theta_0 + \Gamma z$. Under
the Boussinesq approximation, valid when the density contrast across the layer
is small compared with the reference density \\citep{skaar2013boussinesq},
the governing system is
\begin{align}
  \nabla \cdot \mathbf{u} &= 0, \label{eq:continuity} \\
  \frac{\partial \mathbf{u}}{\partial t} + (\mathbf{u} \cdot \nabla)\mathbf{u}
    &= -\nabla p + \nu \nabla^{2} \mathbf{u} + \beta g \theta \, \hat{\mathbf{z}},
    \label{eq:momentum} \\
  \frac{\partial \theta}{\partial t} + (\mathbf{u} \cdot \nabla)\theta
    &= \kappa \nabla^{2} \theta - \Gamma w, \label{eq:scalar}
\end{align}
where $\nu$ is the kinematic viscosity, $\kappa$ the scalar diffusivity, $\beta$
the expansion coefficient, $g$ the gravitational acceleration, and $w$ the
vertical velocity component. The final term in~\eqref{eq:scalar} is the
advection of the background gradient and is the only source of scalar variance
in the configurations studied here.

The buoyancy frequency follows from the background gradient as
\begin{equation}
  N^{2} = \beta g \Gamma,
  \label{eq:buoyancy-frequency}
\end{equation}
and the two dimensionless groups that organise the problem are the friction
Reynolds number $\Rey_\tau = u_\tau h / \nu$ and the bulk Richardson number
\begin{equation}
  \Ri_b = \frac{N^{2} h^{2}}{u_\tau^{2}},
  \label{eq:bulk-richardson}
\end{equation}
with $h$ the channel half height and $u_\tau$ the friction velocity. A local
counterpart, the gradient Richardson number
\begin{equation}
  \Ri_g(z) = \frac{N^{2}}{\left( \partial \langle u \rangle / \partial z \right)^{2}},
  \label{eq:gradient-richardson}
\end{equation}
is the quantity that closures are normally written against, and
Chapter~\ref{chap:results} argues that it is insufficient on its own.

\begin{definition}[Active mixing fraction]
\label{def:active-fraction}
For a domain $\Omega$ and a grid with local spacing $\Delta(\mathbf{x})$, the
active mixing fraction is
\begin{equation*}
  \phi_a = \frac{1}{|\Omega|} \int_{\Omega}
    \mathbb{1}\!\left[ L_O(\mathbf{x}) > 2\Delta(\mathbf{x}) \right] \, d\mathbf{x},
\end{equation*}
the volume fraction over which the Ozmidov scale is resolved by at least two
grid spacings.
\end{definition}

Definition~\ref{def:active-fraction} is the quantity that the refinement
criterion of Section~\ref{sec:method-criterion} is designed to drive towards
unity, and it is reported for every production case in
Table~\ref{tab:case-matrix}.

\section{Discrete kinetic equation}
\label{sec:theory-kinetic}

The momentum field is carried by a population $f_i(\mathbf{x},t)$ defined on a
set of discrete velocities $\mathbf{c}_i$, $i = 0, \dots, 26$, forming the D3Q27
lattice. The evolution equation over a time step $\Delta t$ is
\begin{equation}
  f_i(\mathbf{x} + \mathbf{c}_i \Delta t, t + \Delta t) - f_i(\mathbf{x}, t)
    = \Omega_i^{f}(\mathbf{x},t) + \frac{\Delta t}{2}\left( F_i(t) + F_i(t + \Delta t) \right),
  \label{eq:lbe-momentum}
\end{equation}
where $\Omega_i^{f}$ is the collision operator and $F_i$ the discrete forcing
that carries buoyancy into the momentum equation. The equilibrium population is
the second order Hermite truncation of the Maxwellian,
\begin{equation}
  f_i^{\mathrm{eq}} = w_i \rho \left[
    1 + \frac{\mathbf{c}_i \cdot \mathbf{u}}{c_s^{2}}
      + \frac{(\mathbf{c}_i \cdot \mathbf{u})^{2}}{2 c_s^{4}}
      - \frac{\mathbf{u} \cdot \mathbf{u}}{2 c_s^{2}} \right],
  \label{eq:equilibrium}
\end{equation}
with lattice sound speed $c_s^{2} = \tfrac{1}{3} (\Delta x / \Delta t)^{2}$ and
weights $w_i$ tabulated in Appendix~\ref{app:coefficients}.

The scalar field is carried by a second population $g_j(\mathbf{x},t)$ on the
reduced D3Q7 set, $j = 0, \dots, 6$, which is sufficient because the scalar
equation requires only the first order moment to be recovered correctly. Its
evolution is
\begin{equation}
  g_j(\mathbf{x} + \mathbf{c}_j \Delta t, t + \Delta t) - g_j(\mathbf{x},t)
    = -\frac{1}{\tau_g}\left( g_j - g_j^{\mathrm{eq}} \right)
      - \Delta t \, \Gamma \, w \, \tilde{w}_j,
  \label{eq:lbe-scalar}
\end{equation}
with $g_j^{\mathrm{eq}} = \tilde{w}_j \theta \left[ 1 + \mathbf{c}_j \cdot \mathbf{u} / c_{s,g}^{2} \right]$
and $c_{s,g}^{2} = \tfrac{1}{4}(\Delta x/\Delta t)^2$ for the seven velocity set.

The macroscopic fields are the low order moments
\begin{equation}
  \rho = \sum_i f_i, \qquad
  \rho \mathbf{u} = \sum_i \mathbf{c}_i f_i + \frac{\Delta t}{2} \rho \beta g \theta \hat{\mathbf{z}}, \qquad
  \theta = \sum_j g_j.
  \label{eq:moments}
\end{equation}

\section{Collision in moment space}
\label{sec:theory-mrt}

A single relaxation time operator ties the bulk viscosity, the shear viscosity,
and every higher moment to one rate, and at the Reynolds numbers of interest the
resulting under damping of the ghost moments produces a well documented
checkerboard instability \citep{manov2017mrt}. The momentum population therefore
uses a multiple relaxation time operator. Let $\mathbf{M}$ be the orthogonal
transform that maps the 27 populations to 27 moments $\mathbf{m} = \mathbf{M}\mathbf{f}$,
given explicitly in Appendix~\ref{app:coefficients}. Collision is diagonal in
moment space,
\begin{equation}
  \Omega^{f} = -\mathbf{M}^{-1} \mathbf{S} \left( \mathbf{m} - \mathbf{m}^{\mathrm{eq}} \right),
  \label{eq:mrt-collision}
\end{equation}
where $\mathbf{S} = \mathrm{diag}(s_0, \dots, s_{26})$ collects the relaxation
rates. The rates for the conserved moments are irrelevant and set to unity. The
shear rate $s_\nu$ fixes the viscosity through
\begin{equation}
  \nu = c_s^{2} \left( \frac{1}{s_\nu} - \frac{1}{2} \right) \Delta t,
  \label{eq:viscosity}
\end{equation}
and the scalar relaxation time fixes the diffusivity through the analogous
relation $\kappa = c_{s,g}^{2}(\tau_g - \tfrac{1}{2})\Delta t$. Because the two
populations relax independently, the molecular Prandtl number
$\Pran = \nu / \kappa$ is a free parameter, which is not the case for a single
population formulation with a passive tracer moment.

\begin{proposition}[Second order consistency]
\label{prop:consistency}
Under the scaling $\Delta t \sim \Delta x^{2}$ and for relaxation rates bounded
away from $2$, the Chapman Enskog expansion of~\eqref{eq:lbe-momentum}
and~\eqref{eq:lbe-scalar} recovers~\eqref{eq:continuity} to~\eqref{eq:scalar}
with truncation error $\mathcal{O}(\Delta x^{2})$ and with an error in the
scalar flux that is independent of $s_\nu$.
\end{proposition}

The proof follows the standard multiscale expansion in the Knudsen number
\\citep{iversen2014kinetic, bankole2015chapman}, with
the one non standard step being the treatment of the coupling term. Writing
$f_i = f_i^{\mathrm{eq}} + \epsilon f_i^{(1)} + \epsilon^{2} f_i^{(2)}$ and
collecting orders, the first order momentum balance yields the deviatoric stress
\begin{equation}
  \Pi_{\alpha\beta}^{(1)} = \sum_i c_{i\alpha} c_{i\beta} f_i^{(1)}
    = -\frac{\rho c_s^{2}}{s_\nu}
      \left( \partial_\alpha u_\beta + \partial_\beta u_\alpha \right) \Delta t,
  \label{eq:stress}
\end{equation}
which is the expression that the interface coupling of
Section~\ref{sec:method-coupling} must preserve. The trace free part
of~\eqref{eq:stress} is what carries the physical shear stress, and the
correction to the scalar flux enters only at $\mathcal{O}(\epsilon^{2})$ because
the D3Q7 equilibrium contains no second order velocity term.

\section{Resolution requirement}
\label{sec:theory-resolution}

Combining the Ozmidov scale~\eqref{eq:ozmidov-intro} with the Kolmogorov scale
$\eta = (\nu^{3}/\varepsilon)^{1/4}$ gives the buoyancy Reynolds number
\begin{equation}
  \Rey_b = \left( \frac{L_O}{\eta} \right)^{4/3} = \frac{\varepsilon}{\nu N^{2}},
  \label{eq:buoyancy-reynolds}
\end{equation}
which is the standard measure of whether an inertial range exists between the
two scales. When $\Rey_b$ falls below roughly $20$ the range closes and the flow
becomes intermittent \\citep{ferreira2019intermittency}. The practical consequence for the solver is that the grid
must resolve $L_O$, not $\eta$, in the quiescent regions, and must resolve $\eta$
only in the active sheets. The refinement criterion of
Section~\ref{sec:method-criterion} implements exactly this distinction, and
Table~\ref{tab:case-matrix} reports how much it saves.
